\section{Uniform expected-depth hardness}
\label{sec:uniform-hardness}

We transfer the weighted singleton/pair hardness instance to a uniform
prior by replacing each vertex of weight $w$ with a cloud of $w$
hypotheses. Private queries can occur anywhere in an adaptive tree,
so the key step is to move them off the distinguished spine while
preserving or reducing cost. This normalization exposes whole clouds
as matching groups and yields an exact identity: the cloud optimum is
the weighted optimum plus a computable internal-cloud cost. Here raw
uniform cost denotes the sum of leaf depths, before division by the
number of hypotheses.

\subsection{Exact cloud replacement: normalization before accounting}

\subsubsection{The construction}

Consider the weighted singleton/pair linear forest instance of
Lemma~\ref{lem:pair-invoice}, with vertices $0,1,\ldots,N$ and
positive integer weights $w_0,\ldots,w_N$. Vertex $0$ is
distinguished. Order the nonzero vertices along each path and
component by component, so every edge joins consecutive indices.
Here a cloud $C_i$ is the block of $w_i$ hypotheses representing
vertex $i$, its anchor is the point with offset $j=0$, and a private
query is one of the supplied lines that isolates a single nonanchor
point. The singleton and pair queries represent, respectively, a
vertex and an edge of the weighted forest. The invoice notation
$J(\mathcal M)$ and $\mathcal C(J(\mathcal M))$ is as in
Lemma~\ref{lem:pair-invoice}; it records the weighted instance left
after choosing the matching $\mathcal M$ of pair queries.
Choose the cloud spacing $Q$ by the final calibration below and define
\[
 n=\sum_{i=0}^N w_i,\qquad
 C_i=\{(Qi+j,(Qi+j)^2):0\le j<w_i\}.
\]
The point with $j=0$ is the anchor of cloud $C_i$.
For every nonzero vertex supply its whole-cloud singleton query;
for every linear forest edge supply its whole-cloud pair query.
For consecutive cloud indices $a,\ldots,b$, these are the negative
answers of
\begin{equation}\label{eq:uniform-cloud-membership}
 4y-4Q(a+b)x+Q^2(2a-1)(2b+1)=0.
\end{equation}
The value on the parabola factors as
$4(x-Q(a-\tfrac12))(x-Q(b+\tfrac12))$.
The product is negative exactly between its two roots. Cloud $c$ has
abscissae in $[Qc,Qc+w_c-1]$, with $w_c-1<Q/4$; hence every point
in clouds $a$ through $b$ lies strictly between the roots, while each
other cloud lies strictly outside. Since $a>0$, cloud $0$ is
excluded. Thus the negative answer selects exactly the designated
cloud or adjacent pair, and no point lies on the query line.

For every nonanchor point of abscissa $t$, including those of
cloud $0$, supply the private line
\[
 2y-4tx+2t^2-1=0.
\]
Its value on $(x,x^2)$ is $2(x-t)^2-1$; it is negative at $x=t$
and positive at every other integer abscissa, since then $|x-t|\ge1$.
It therefore isolates precisely that point, with no point on the line.
No private line is supplied for an anchor. All answers are strict
and binary. The whole-cloud singleton queries distinguish each
nonzero anchor from all other anchors and cloud $0$; the private
queries then separate every nonanchor from the anchor in its cloud.
Consequently the query family identifies every hypothesis. There are
$n=\sum_iw_i$ hypotheses, so the uniform prior gives each weight $1/n$.

Put
\begin{equation}\label{eq:uniform-cloud-offset}
 D(w)=\frac{(w-1)(w+2)}2,\qquad
 \Delta=\sum_{i=0}^N D(w_i).
\end{equation}
Algorithm~\ref{alg:uniform-cloud-lift} gives the construction and
the threshold conversion.

\begin{algorithm}[H]
\caption{Uniform cloud replacement of a weighted linear forest instance}
\label{alg:uniform-cloud-lift}
\KwIn{Vertices $0,\ldots,N$ with each path in order, component by component,
positive integer weights $w_i$, and a raw weighted threshold $k$}
\KwOut{Supplied-line instance with uniform prior and threshold
$(k+\Delta)/n$}
$Q\gets4\max_i w_i$; $n\gets\sum_iw_i$;
$\HH\gets\varnothing$; $\LL\gets\varnothing$\;
\For{$i=0,\ldots,N$}{
  \For{$j=0,\ldots,w_i-1$}{
    $t\gets Qi+j$; add $(t,t^2)$ to $\HH$\;
    \If{$j>0$}{
      Add $2y-4tx+2t^2-1=0$ to $\LL$\;
    }
  }
}
For each $i>0$, add \eqref{eq:uniform-cloud-membership}
with $(a,b)=(i,i)$ to $\LL$\;
For each linear forest edge $\{i,i+1\}$, add
\eqref{eq:uniform-cloud-membership} with $(a,b)=(i,i+1)$ to $\LL$\;
$\Delta\gets\sum_{i=0}^N(w_i-1)(w_i+2)/2$\;
\KwRet $(\HH,\LL)$, prior $1/n$, and threshold $(k+\Delta)/n$\;
\end{algorithm}

\begin{theorem}[Exact uniform cloud transfer]\label{thm:uniform-cloud-transfer}
Let $\OPT_{\mathrm{raw}}^{\mathrm{weighted}}$ be the minimum
raw weighted cost of the singleton/pair linear forest instance.
For its uniform cloud replacement,
\begin{equation}\label{eq:uniform-cloud-transfer}
 n\OPT_{\mathrm E}^{\mathrm{cloud}}
   =\OPT_{\mathrm{raw}}^{\mathrm{weighted}}+\Delta.
\end{equation}
Equivalently, if $J(\mathcal M)$ has the pair and singleton masses
from Lemma~\ref{lem:pair-invoice}, the raw uniform optimum is
\[
 \Delta+\min_{\mathcal M}
 \left\{w_0(N-|\mathcal M|)
   +\mathcal C(J(\mathcal M))
   +\sum_{\{u,v\}\in\mathcal M}(w_u+w_v)\right\}.
\]
\end{theorem}

\begin{proofstrategy}
\begin{enumerate}
\item Rotate a private singleton past the next query. Equal point
weights make its one-unit depth increase no larger than the savings
on the opposite, nonempty answer branch.
\item Push all external private queries off the distinguished anchor's
spine and move its own private chain to the end. Every first-hit cloud
is then intact, so the first-hit pairs form a matching.
\item Sum exact one- and two-cloud local costs, then sort group sizes.
The result is the weighted matching cost plus the fixed offset
$\Delta$; the same construction attains it.
\end{enumerate}
\end{proofstrategy}

\subsubsection{Three supporting statements}

Uniform weights make one local exchange cost nonincreasing: delaying
an isolated point by one query saves a query for a nonempty opposite
branch. Since every point has the same weight, the exchange changes
the raw cost only through the depths immediately above the attached
subtrees; their internal costs cancel. This is the tool for
normalizing arbitrary trees.

\begin{lemma}[Uniform singleton rotation]\label{lem:uniform-singleton-rotation}
Suppose a proper identifying tree on a uniformly weighted state $S$
first asks a supplied query isolating $x\in S$, and the other child
first asks a query $r$ that is proper on $S\setminus\{x\}$.
These two queries can be exchanged, placing the singleton query
in the answer branch of $r$ containing $x$, without increasing
the raw cost.
\end{lemma}
Once clouds reach their first hit intact, their internal cost can
be accounted for exactly. We need both the single-cloud value and
a check that early private queries cannot improve a two-cloud exit.

\begin{lemma}[Local uniform cloud costs]\label{lem:uniform-cloud-local}
The optimal raw local cost on one intact cloud of size $w$ is
$D(w)$. On two intact clouds of sizes $u,v$, it is
\[
 D(u)+D(v)+u+v.
\]
\end{lemma}

The rotations must establish a global normal form before costs
are summed. External private queries leave the distinguished spine,
while its own private chain moves to the end, exposing whole-cloud
matching groups.

\begin{proposition}[A normalized matching spine]
\label{prop:uniform-cloud-spine}
An optimal uniform cloud tree can be chosen so that its distinguished
anchor's path first asks only whole-cloud queries, with every newly
hit cloud intact, and then asks only cloud-$0$ private queries.
The first-hit pairs form a matching $\mathcal M$. For this matching,
the best raw cost is
\[
 \Delta+w_0(N-|\mathcal M|)+\mathcal C(J(\mathcal M))
          +\sum_{\{u,v\}\in\mathcal M}(w_u+w_v).
\]
Every matching has a tree attaining this invoice.
\end{proposition}

\begin{proof}[Proof of Theorem~\ref{thm:uniform-cloud-transfer}]
Lemma~\ref{lem:uniform-singleton-rotation} permits the normalization
in Proposition~\ref{prop:uniform-cloud-spine} without increasing cost.
Lemma~\ref{lem:uniform-cloud-local} supplies its exact local offset.
Minimize the resulting invoice over matchings and apply
Lemma~\ref{lem:pair-invoice}. This gives
$\OPT_{\mathrm{raw}}^{\mathrm{cloud}}
=\OPT_{\mathrm{raw}}^{\mathrm{weighted}}+\Delta$.
Dividing by $n$ proves the displayed expected-depth identity.
\end{proof}

\subsubsection{Technical proofs of the supporting statements}

\begin{proofstrategy}
First, compare the two possible query orders by their per-state depth
changes; uniform weights make the singleton rotation nonincreasing.
Next, compute the forced private-chain cost for one cloud and use
induction to show that a whole-cloud separator is optimal for two.
Finally, apply those rotations to normalize the distinguished path,
then sum each first-hit group's local cost and order the groups by
decreasing mass. This yields the matching cost formula and its attaining
construction.
\end{proofstrategy}

\begin{proof}[Proof of Lemma~\ref{lem:uniform-singleton-rotation}]
Let $A,B$ be the two nonempty answer classes of $r$ on
$S\setminus\{x\}$, and choose their names so that $x$ has the
answer of $A$. In the original tree, $x$ has relative depth one
and every hypothesis in $A\cup B$ passes through both queries.
Ask $r$ first, put the private singleton query in its
$A\cup\{x\}$ branch, and attach the original subtrees on $A$ and
$B$. Originally $x$ has relative depth one and points of $A$ and $B$
have relative depth two before their attached subtrees. Afterwards
$x$ has relative depth two, the points of $A$ still have depth two,
and the points of $B$ have depth one before their unchanged subtrees.
Thus the raw cost change is
\[
 (2-1)+|A|(2-2)+|B|(1-2)=1-|B|\le0,
\]
because $B$ is nonempty. The resulting queries are
proper and are members of the original supplied family.
\end{proof}

If the opposite class $B$ has one point, the rotation has zero cost;
if it has more than one, the cost decreases. The lemma concerns
equal weights, so we use the rotation only when every point contributes
the same unit cost at each level.

\begin{proof}[Proof of Lemma~\ref{lem:uniform-cloud-local}]
We first compute the forced private-chain cost within one cloud.
For two clouds, we then compare separating them immediately with
every possible private-first choice; this verifies the offset without
assuming a local normal form.

\begin{claim}[A single cloud forces a private chain]
\label{clm:uniform-cloud-chain}
Its raw local cost is $D(w)$, including zero cost when $w=1$.
\end{claim}
\begin{proof}
Within one cloud, every whole-cloud query is constant, and every
nonconstant query isolates one of its nonanchors. The only possible
proper code is a chain: every private query has a singleton outcome,
and the other outcome retains the anchor and all points not yet
isolated. Thus the nonanchors exit at depths $1,\ldots,w-1$, and
the anchor exits at depth $w-1$. The sum of depths is
$1+2+\cdots+(w-1)+(w-1)$. For $w\ge2$,
\[
 \sum_{j=1}^{w-1}j+(w-1)
 =\frac{w(w-1)}2+(w-1)
 =\frac{(w-1)(w+2)}2=D(w).
\]
For $w=1$ there is no query and both the sum and $D(1)$ are zero.

\end{proof}
Separating the clouds gives the proposed cost directly. The remaining
question is whether isolating a private point first could save more;
induction accounts for the delay imposed on the other cloud.

\begin{claim}[A whole-cloud separator is optimal on two clouds]
\label{clm:uniform-cloud-separator}
Private-first choices cannot beat $D(u)+D(v)+u+v$.
\end{claim}
\begin{proof}
On two clouds a nonconstant whole-cloud query separates the
clouds. Its root query is traversed by all $u+v$ points, after which
the clouds cost $D(u)+D(v)$; this attains the displayed formula.
Every other nonconstant query isolates one nonanchor. Induct on
$u+v$, with the whole-cloud separator available at every induction
state.
If a private point of the $u$ cloud is queried first, the
remaining two-cloud state has sizes $u-1,v$ and still has its
anchors and whole-cloud separator. Using induction, this choice
costs
\[
 u+v+D(u-1)+D(v)+(u-1)+v.
\]
A private point exists in this cloud only if $u\ge2$. The
single-cloud formula of Claim~\ref{clm:uniform-cloud-chain} gives
\[
 D(u)-D(u-1)
 =\frac{(u-1)(u+2)-(u-2)(u+1)}2=u.
\]
Subtracting the proposed separator cost from the private-first cost
therefore gives
\begin{align*}
 &[u+v+D(u-1)+D(v)+(u-1)+v]
       -[D(u)+D(v)+u+v]\\
 &\qquad=D(u-1)-D(u)+(u-1)+v
          =-u+(u-1)+v=v-1\ge0.
\end{align*}
The second equality substitutes $D(u)-D(u-1)=u$, and the final
inequality uses the presence of the other cloud, namely $v\ge1$.
Thus any saving on the queried cloud is paid for by delaying the
other cloud's points.
For a private point in the other cloud, the identical calculation
with $u,v$ interchanged gives $u-1\ge0$. The base case $u=v=1$ has raw cost two.
When one size is one it has no private point, and the same
induction applies to the other cloud. Queries selecting other
clouds are constant or separate these two whole clouds, so the
argument covers the entire supplied family.
\end{proof}
\end{proof}


\begin{proof}[Proof of Proposition~\ref{prop:uniform-cloud-spine}]
The two normalization steps have different destinations. A nonzero
cloud's private query should leave the distinguished path at its
whole-cloud hit; a cloud-$0$ private query stays on that path but
moves behind its whole-cloud queries.

\begin{claim}[External private queries can be pushed off the spine]
\label{clm:uniform-cloud-external-private}
No nonzero cloud needs to lose a private point before its first whole-cloud hit.
\end{claim}
\begin{proof}
Start with a proper optimal uniform tree and follow the anchor
of cloud $0$. Suppose a private query for a nonanchor $x$ of a
nonzero cloud occurs on this path. Its own cloud anchor is still
on the path: $x$ agrees with its anchor on every whole-cloud query,
and private queries never remove an anchor.
At each step, let $r$ be the next query on the anchor-$0$ path.
The private query isolates $x$, so the anchor-$0$ path continues in
the other child; the state reaching $r$ is precisely that non-$x$
answer class. Since $r$ is proper on $S\setminus\{x\}$, it remains
proper on this state. Apply
Lemma~\ref{lem:uniform-singleton-rotation} to move the private query
for $x$ past $r$. This also applies when $r$ is private: it isolates
a different point, and the same depth comparison is nonincreasing.
Eventually a whole-cloud query separates $x$'s
anchor from the anchor of cloud $0$: otherwise the tree would
not identify both anchors. The point $x$ has the answer of its
own anchor to that query. After the rotation its private query
is in the hit branch, off the anchor-$0$ path.

Each completed operation removes one external private query from
the path and introduces none. Repeating over the finite set of such
queries terminates; each individual push crosses only the finite
path prefix up to the separating whole-cloud query.
\end{proof}
External private queries have now been removed from the spine.
The remaining private points agree with the distinguished anchor on
every whole-cloud query, so the same rotation postpones their chain
without moving it to an exit.

\begin{claim}[The distinguished private chain can be moved to the end]
\label{clm:uniform-cloud-distinguished-private}
All whole-cloud queries can precede cloud $0$'s private queries.
\end{claim}
\begin{proof}
Next move each cloud-$0$ private query past the next whole-cloud
query on the anchor path by the rotation lemma. Its point always has
the answer of the anchor of cloud $0$ to that query, so the path
continues through the whole-cloud query and the private query remains
on the path behind it. This second process terminates because each
rotation decreases by one the finite number of private/whole-cloud
inversions (pairs in which a private query precedes a whole-cloud
query on this path).
Every rotation is cost nonincreasing by
Lemma~\ref{lem:uniform-singleton-rotation}; together with
Claim~\ref{clm:uniform-cloud-external-private}, this establishes the
required whole-cloud-first spine.

\end{proof}
With private queries in these positions, each first-hit cloud is
intact. We can therefore use the exact local costs and sort group
masses; constructing the same arrangement proves that the resulting
invoice is attainable.

\begin{claim}[The normalized groups attain the matching cost]
\label{clm:uniform-cloud-invoice}
The minimum normalized raw cost is the expression in Theorem~\ref{thm:uniform-cloud-transfer}.
\end{claim}
\begin{proof}
Claims~\ref{clm:uniform-cloud-external-private} and
\ref{clm:uniform-cloud-distinguished-private} imply that all newly
hit clouds are intact at their first
whole-cloud query: no private query for any of their points
has occurred on the path earlier. First-hit groups are whole
singleton clouds or whole adjacent pairs, and the pairs form
a matching $\mathcal M$: the only supplied whole-cloud queries are
singleton queries and pair queries on forest edges, and a cloud can
belong to only one first-hit group. A query first hitting just one
cloud (for example, a pair query whose other endpoint has already
left the path) can be replaced by its supplied whole-cloud singleton query,
which has the same split on the current state. There are $|\mathcal M|$ pair groups and
$N-2|\mathcal M|$ singleton groups, hence
$k=N-|\mathcal M|$ whole-cloud spine queries. Every point of cloud
$0$ traverses all of them, contributing $w_0k$, followed by its
private-chain cost $D(w_0)$. A singleton cloud contributes $D(w_i)$
after its first hit. By Lemma~\ref{lem:uniform-cloud-local}, a pair
contributes $D(w_u)+D(w_v)+w_u+w_v$ after its first hit. Summing
these local costs over disjoint groups therefore gives
$\Delta+\sum_{\{u,v\}\in\mathcal M}(w_u+w_v)$.

If the group's total masses in spine order are $a_1,\ldots,a_k$,
their first-hit depths contribute $\sum_j j a_j$: every point in
group $j$ passes through exactly the first $j$ whole-cloud queries
before it exits the distinguished path. For consecutive
$a<b$, exchanging their positions changes their contribution from
$ja+(j+1)b$ to $jb+(j+1)a$, a decrease of $b-a>0$.
Thus their minimum is $\mathcal C(J(\mathcal M))$ in nonincreasing
mass order. Adding the witness, first-hit, and local contributions
gives the stated matching cost as a lower bound. Every
$D(w_i)$ occurs once because the first-hit groups partition the
nonzero clouds and the final private chain accounts for cloud $0$;
this is precisely why their sum is the matching-independent offset
$\Delta$ of \eqref{eq:uniform-cloud-offset}.

Conversely, for any matching, query its disjoint pairs and
unmatched singleton clouds in decreasing total cloud size.
Each selected pair query or singleton query peels exactly its group
from the still-unhit clouds on the distinguished path; the branch
that exits the path contains exactly that group's cloud or clouds.
Attach the
optimal local codes of Lemma~\ref{lem:uniform-cloud-local} in those
branches and identify cloud $0$ by its private chain at the end. This
realizes the decreasing-mass first-hit order and attains the matching
invoice plus $\Delta$. Taking the minimum proves equality.
\end{proof}
\end{proof}

\subsection{Uniform-prior strong NP-completeness}

\begin{theorem}[Strong NP-completeness under a uniform prior]
\label{thm:strong-uniform-average}
Deciding whether the optimal expected depth is at most a rational
threshold is strongly NP-complete under a uniform prior, even for
integer-coordinate points on a parabola, polynomially bounded
integer line coefficients, and $O(n)$ supplied strict binary lines.
\end{theorem}

\begin{proofstrategy}
\begin{enumerate}
\item Apply the exact cloud replacement to the hard weighted forest
instances of Theorem~\ref{thm:strong-average}.
\item Their polynomial integer weights allow explicit replication in
polynomial time. Separated parabola clouds give strict binary queries
with polynomial coordinates and coefficients.
\item Convert the raw threshold to $(k+\Delta)/n$ by the transfer
identity. A proper tree provides a polynomial certificate.
\end{enumerate}
\end{proofstrategy}

\subsubsection{Supporting size and encoding propositions}

Replacing a weight by that many points is polynomial only when
the weights themselves are polynomially bounded. The hard weighted
family satisfies this stronger requirement, so explicit clouds are feasible.

\begin{proposition}[Polynomial replication]
\label{prop:uniform-replication}
For a hard weighted instance with $n_0=N+1$ hypotheses, the explicit
cloud replacement has $n=O(n_0^5)$ points and is constructible in
polynomial time on the restricted 3-Partition source instances.
\end{proposition}

Replication increases the number of points, and spacing the clouds
increases their coordinates. We must bound both effects while retaining
strict binary answers and the supplied-query geometry.

\begin{proposition}[Strong geometric encoding]
\label{prop:uniform-geometric-encoding}
The replacement has $O(n)$ supplied strict binary lines. Its integer
abscissae are $O(n^2)$; its ordinates and line coefficients are $O(n^4)$.
The points remain in strictly convex position on a parabola.
\end{proposition}

The transfer identity adds a fixed cost to every matching. Adding
the same offset to the threshold preserves the decision answer;
polynomial threshold data and tree certificates complete the hardness
argument.

\begin{proposition}[Threshold conversion and NP certificates]
\label{prop:uniform-threshold-certificate}
For the hard weighted instances and thresholds of
Theorem~\ref{thm:strong-average}, the rational threshold
$(k+\Delta)/n$ has polynomially bounded
numerator and denominator. Its yes/no answer agrees with that of the
weighted raw threshold $k$ for arbitrary adaptive trees, and the
uniform decision problem belongs to NP.
\end{proposition}

\begin{proof}[Proof of Theorem~\ref{thm:strong-uniform-average}]
Apply Algorithm~\ref{alg:uniform-cloud-lift} to the weighted hard
instances. Proposition~\ref{prop:uniform-threshold-certificate} and
Theorem~\ref{thm:uniform-cloud-transfer} preserve the yes/no answer.
Propositions~\ref{prop:uniform-replication}
and~\ref{prop:uniform-geometric-encoding} give a polynomial reduction
with polynomial numerical data. The certificate proposition supplies
NP membership, proving strong NP-completeness.
\end{proof}

\subsubsection{Technical proofs of the encoding propositions}

\begin{proofstrategy}
The replication bound follows by summing the calibrated weights.
For the geometric bound, count the queries and substitute the spacing
bound into the point and line formulas. For threshold size and NP
membership, use the transfer identity, bound $\Delta$ by $n^2$, and
encode a proper identifying tree by its $n-1$ internal nodes.
\end{proofstrategy}

\begin{proof}[Proof of Proposition~\ref{prop:uniform-replication}]
Proposition~\ref{prop:average-internal-encoding} and its calibration
give nonzero weights $O(n_0^3)$ and
distinguished weight $O(n_0^5)$. There are $N=n_0-1$ nonzero
weights, so their sum is $O(n_0\cdot n_0^3)=O(n_0^4)$.
Consequently
$n=w_0+\sum_{i>0}w_i=O(n_0^5)+O(n_0^4)=O(n_0^5)$. The restricted source makes $n_0$ polynomial
in its combinatorial size, so all points and queries can be listed
explicitly in polynomial time.
\end{proof}

\begin{proof}[Proof of Proposition~\ref{prop:uniform-geometric-encoding}]
There are $n-(N+1)$ private lines and at most $2N$ whole-cloud
lines, hence $O(n)$ supplied queries. The calibrated spacing bounds
all coordinates and coefficients as verified in the calibration
below. On every private line the value is $2(x-t)^2-1$, which is
$-1$ at its target and at least $1$ at every other integer
abscissa. The whole-cloud lines have the factored values in
the construction, with roots strictly between clouds, so their
answers are also strict. Each nonzero cloud is separated from the
rest by its singleton query, and each nonanchor is separated from
every other point by its private query; hence all hypotheses are
distinguishable. Finally, for any point $(t,t^2)$, the tangent
$y=2tx-t^2$ satisfies $s^2-(2ts-t^2)=(s-t)^2>0$ at every other
abscissa $s$. Thus each sampled point is an exposed vertex of the
convex hull, proving strict convex position.
\end{proof}

\begin{proof}[Proof of Proposition~\ref{prop:uniform-threshold-certificate}]
Theorem~\ref{thm:uniform-cloud-transfer} gives
$\OPT_{\mathrm{raw}}^{\mathrm{weighted}}\le k$
if and only if
$\OPT_{\mathrm E}^{\mathrm{cloud}}\le(k+\Delta)/n$.
For the weighted instance, Proposition~\ref{prop:average-internal-encoding}
gives $k=O(n_0^6)$ and total weight $\Omega=O(n_0^5)$, where
$n_0=N+1$. The cloud construction has $n\ge n_0$, since every weight
is a positive integer, so $k=O(n^6)$. For $w\ge1$,
\[
 w^2-D(w)=\frac{w^2-w+2}{2}\ge0.
\]
Indeed $w^2-w=w(w-1)\ge0$, so each cloud's offset is at most
the square of its size.
Thus
\[
 \Delta\le\sum_iw_i^2
 \le\left(\sum_iw_i\right)^2=n^2,
\]
where the second inequality adds the nonnegative cross products.
Hence the unreduced numerator satisfies $k+\Delta=O(n^6)$, while
the denominator is $n$. Both integers have $O(\log n)$ binary digits,
and reducing the fraction can only decrease their magnitudes.
A proper identifying tree has $n-1$ internal nodes and is a polynomial
certificate. Evaluate each supplied answer and sum the leaf depths
to verify its rational expected cost in polynomial time.
\end{proof}

\subsubsection{Calibration of the cloud spacing}

\begin{parameterchoice}
Choose $Q=4\max_iw_i$. Every cloud has width less than $Q/4$;
the half-anchor boundaries of \eqref{eq:uniform-cloud-membership}
therefore miss all points and select exactly the stated clouds.
Since $Q\le4n$ and $N\le n$, every abscissa is $O(n^2)$.
Squaring it bounds the ordinates by $O(n^4)$, and the formulas for
both query types bound all integer coefficients by $O(n^4)$.
\end{parameterchoice}

\begin{remark}[Limits of the reduction]
The cloud identity adds the same nonnegative internal cost $\Delta$
to every matching. Let $T_w$ be the matching returned by the weighted
EPTAS and let $T_u$ be its lifted tree on the uniform cloud instance.
The exact cost identity gives
\[
 c(T_u)=c(T_w)+\Delta
 \le(1+\varepsilon)\OPT_{\rm weighted}+\Delta
 \le(1+\varepsilon)(\OPT_{\rm weighted}+\Delta)
 =(1+\varepsilon)\OPT_{\rm uniform}.
\]
Thus the lift preserves the EPTAS as well as exact hardness. A
constant-factor inapproximability result would need a different
construction.
\end{remark}
